\section{Edit deficits, ordered rows and binary transfer}
\label{sec:preliminaries}

Both trees need four elementary tools. Subsequence deficits and an
alignment lemma turn separation of child blocks into separation of
whole words, even when matches cross row boundaries. A cut
representation permits the later Fourier estimates, and a prime count
supplies the periods. Finally, a short binary code transfers either
tree's obstruction to binary words with a controlled width loss.

\subsection{Subsequence matchings}

A common subsequence matching of words $x=x_1\cdots x_a$ and
$y=y_1\cdots y_b$ is a list of pairs
$(i_1,j_1),\ldots,(i_s,j_s)$ whose two index lists are strictly
increasing and for which $x_{i_t}=y_{j_t}$ for every $t$. Its largest
possible size is $\LCS(x,y)$. The insertion--deletion distance is
\[
 \Delta(x,y)=|x|+|y|-2\LCS(x,y).
\]
Indeed, symbols of $x$ that survive an insertion--deletion script form
a common subsequence; conversely, retaining a longest common subsequence
and deleting and inserting the remaining symbols attains this value.
Replacing each substitution by one deletion and one insertion gives
\begin{equation}\label{eq:metric-conversion}
 \ED(x,y)\leq\Delta(x,y)\leq2\ED(x,y).
\end{equation}
For words of equal length $n$, put $r_n(x,y)=n-\LCS(x,y)$. Thus
\begin{equation}\label{eq:deficit}
 r_n(x,y)\leq\ED(x,y)\leq2r_n(x,y).
\end{equation}
A matching of size $n-E$ between two length-$n$ words leaves exactly
$E$ unmatched positions on each side.

\subsection{The common alignment lemma}

Suppose each row consists of one block from each of several pairwise
disjoint alphabets, in the same order. An equal-symbol match must then
join blocks with the same index, although their row indices can differ.
Products of words in the next lemma mean concatenation in increasing
index order.

\Needspace{15\baselineskip}
\begin{lemma}[Rows with ordered tags]\label{lem:tagged}
Let $T,B\geq1$ and $m\geq2$ be integers, and let
$\Gamma_1,\ldots,\Gamma_m$ be pairwise disjoint alphabets. For
$1\leq a,b\leq T$ and $1\leq i\leq m$, take
$X_{a,i},Y_{b,i}\in\Gamma_i^B$, and define
\[
 X=\prod_{a=1}^T(X_{a,1}\cdots X_{a,m}),\qquad
 Y=\prod_{b=1}^T(Y_{b,1}\cdots Y_{b,m}),\qquad N=TmB.
\]
Suppose $d\geq0$ is real and
\begin{equation}\label{eq:local-tagged}
 \sum_{i=1}^m r_B(X_{a,i},Y_{b,i})\geq d
 \qquad(1\leq a,b\leq T).
\end{equation}
With $g=(m-1)B$, one has
\begin{equation}\label{eq:tagged-result}
 r_N(X,Y)\geq\frac{Tdg}{d+g}
 \geq\frac T2\min(d,g).
\end{equation}
In particular, $d\leq g$ implies $r_N(X,Y)\geq Td/2$.
\end{lemma}

\begin{proof}
Fix an increasing matching and let $E$ be its unmatched count on either
side. Call a source row split if it has matches in at least two target
rows, and let $s$ be the number of split source rows.

A source row that is not split and has matches sends them all into one
target row, say $b$. The disjoint alphabets confine matches from its
block $i$ to block $i$ in that target row. The restriction of the
matching to those two blocks has size at most their $\LCS$, even if
other source rows also use the target block. Condition
\eqref{eq:local-tagged} therefore leaves at least $d$ unmatched source
positions in this row. A source row with no matches loses all $mB$
positions, again at least $d$, since \eqref{eq:local-tagged} implies
$d\leq mB$. The nonsplit rows are disjoint, so
\begin{equation}\label{eq:tagged-unsplit}
 E\geq d(T-s).
\end{equation}

For a split source row, list its matches in order and choose two
consecutive entries at a transition between target rows. They are also
consecutive in the entire matching: every source position between them
lies in the same source row, so an intervening global match would have
appeared in its list. Let $i,i'$ be their source block indices. Their
order gives $i\leq i'$, and disjoint alphabets give the same two block
indices in the target. Between their target positions lie at least
\[
 (m-i+i'-1)B\geq(m-1)B=g
\]
positions in complete blocks. This is the count for adjacent target
rows; intervening rows only increase it. All positions in that open
target interval are unmatched, because the selected matches are
globally consecutive. The intervals selected for different source rows
are disjoint by strict monotonicity of the matching. Hence
\begin{equation}\label{eq:tagged-split}
 E\geq sg.
\end{equation}

Since $g>0$, \eqref{eq:tagged-split} gives $s\leq E/g$. Substitute this
in \eqref{eq:tagged-unsplit} to obtain
$E\geq d(T-E/g)$, and rearrange:
$E\geq Tdg/(d+g)$. If $d\leq g$, then $dg/(d+g)\geq d/2$; if
$g\leq d$, it is at least $g/2$. These bounds hold for every matching,
including when $d=0$ or $T=1$, and therefore for a longest one.
\end{proof}

The two inequalities in the proof bound the same number $E$ using
different words: nonsplit rows give source losses, and split rows give
target losses. Equal total lengths make the unmatched counts equal.
The fixed order, equal block lengths and disjoint tags supply the
target gaps; no alignment of row boundaries is assumed.
Figure~\ref{fig:tagged-gap} illustrates the gap created by a split row.

\begin{figure}[ht]
\centering
\begin{tikzpicture}[x=0.98cm,y=0.85cm,>=Latex,
 every node/.style={font=\small}]
 \foreach \block in {1,...,5} {
  \draw ({\block+1.5},2) rectangle ({\block+2.5},2.6);
  \node at ({\block+2},2.3) {$\block$};
 }
 \node[anchor=east] at (2.3,2.3) {source row $a$};
 \foreach \block in {1,...,5} {
  \draw ({\block-1},0) rectangle (\block,0.6);
  \draw ({\block+4},0) rectangle ({\block+5},0.6);
 }
 \foreach \block in {3,4,5} {
  \fill[blue!15] ({\block-1},0) rectangle (\block,0.6);
 }
 \foreach \block in {1,2,3} {
  \fill[blue!15] ({\block+4},0) rectangle ({\block+5},0.6);
 }
 \foreach \block in {1,...,5} {
  \node at ({\block-0.5},0.3) {$\block$};
  \node at ({\block+4.5},0.3) {$\block$};
 }
 \draw[very thick] (5,0) -- (5,0.6);
 \node at (2.5,-0.4) {target row $b$};
 \node at (7.5,-0.4) {target row $b+1$};
 \draw[blue,thick,->] (4,2) -- (1.5,0.6);
 \draw[blue,thick,->] (6,2) -- (8.5,0.6);
 \fill[blue] (4,2) circle (1.5pt);
 \fill[blue] (6,2) circle (1.5pt);
 \draw[blue,<->] (1.5,0.95) -- (8.5,0.95);
 \node[fill=white,inner sep=2pt] at (5,0.95) {unmatched target gap};
\end{tikzpicture}
\caption{A split source row in Lemma~\ref{lem:tagged}, with $m=5$.
The two arrows are consecutive matches from source blocks $i=2$ and
$i'=4$. Their target positions lie in different rows, so all positions
between them are unmatched. The six shaded complete blocks contribute
$(m-i+i'-1)B=6B\geq(m-1)B$ positions. The diagram is schematic; target
rows further apart would only enlarge the gap.}
\label{fig:tagged-gap}
\end{figure}


\subsection{Cuts and prime counts}

We record two background facts for both constructions. The first turns
an $\ellone$ distance on a finite set into a nonnegative sum of binary
separation tests. This standard cut description is stated, for example,
in Section~4 of Linial, London and Rabinovich
\cite{LinialLondonRabinovich1995}; the proof below includes a summability
bound for maps with infinitely many coordinates.

\begin{lemma}[Finite cut representation]\label{lem:cuts}
Let $S$ be finite and $f:S\to\ellone$. There are nonnegative finite
weights $a_C$, indexed by nonempty proper subsets $C\subset S$, such that
\[
 \|f(x)-f(y)\|_1
 =\sum_{\varnothing\ne C\subsetneq S}
 a_C|\one_C(x)-\one_C(y)|.
\]
Moreover,
\[
 \sum_{\varnothing\ne C\subsetneq S}a_C
 \leq\sum_{\{x,y\}\subset S,\ x\ne y}\|f(x)-f(y)\|_1,
\]
where the sum on the right is over unordered pairs.
\end{lemma}
\begin{proof}
For each real coordinate $f_j$,
\[
 |f_j(x)-f_j(y)|
 =\int_{\R}|\one_{\{f_j>t\}}(x)-\one_{\{f_j>t\}}(y)|\,dt.
\]
Discard empty and full level sets, which contribute zero, then sum over
coordinates and group the remaining level sets by their subset of $S$.
All terms are nonnegative, so monotone convergence permits this
grouping. Every nontrivial cut separates at least one unordered pair.
Summing the identity over those pairs bounds the total weight of the
nontrivial cuts by the stated finite sum.
\end{proof}

The second fact supplies all primes used below. We use the elementary
binomial-coefficient argument, following the classical prime-product
method found, for example, in Erd\H{o}s \cite{Erdos1932}. Write $\pi(x)$
for the number of primes at most $x$.

\begin{lemma}[Prime supply]\label{lem:primes}
There is an absolute $c>0$ such that
$\pi(x)\geq cx/\log x$ for all sufficiently large real $x$.
Consequently, for all sufficiently large integers $k$, the interval
$[k^4,k^5]$ contains at least $20k^2$ primes and $[k^2,k^3]$ contains
more than $1600k$ primes.
\end{lemma}
\begin{proof}
For an integer $a\geq1$, the largest coefficient in the expansion of
$(1+1)^{2a}$ is the central one, so
$\binom{2a}{a}\geq4^a/(2a+1)$. The exponent of a prime $p$ in this
coefficient is
\[
 \sum_{s\geq1}
 \left(\left\lfloor\frac{2a}{p^s}\right\rfloor
 -2\left\lfloor\frac a{p^s}\right\rfloor\right).
\]
Each summand is zero or one and vanishes for $p^s>2a$. Hence the whole
$p$-power dividing $\binom{2a}{a}$ is at most $2a$, and
$\binom{2a}{a}\leq(2a)^{\pi(2a)}$. Taking logarithms gives
$\pi(2a)\geq c(2a)/\log(2a)$ for large $a$. Taking
$a=\lfloor x/2\rfloor$ yields the asserted lower bound for real $x$.
Finally subtract $\pi(k^4)\leq k^4$ from the lower bound at $k^5$,
and $\pi(k^2)\leq k^2$ from the lower bound at $k^3$. The respective
differences eventually exceed $20k^2$ and $1600k$.
\end{proof}

\subsection{A short binary code}

The trees initially use leaf labels and hence a growing alphabet. The
following code will suffice for both binary lower bounds. It preserves
subsequence deficits from below, while its expansion is at most its
block width. For the tree alphabets that width will be $O(k\log k)$,
small compared with their eventual distortion lower bounds.

\begin{lemma}[Delimiter transfer]\label{lem:delimiter}
Let $\Gamma$ be a finite alphabet of size $A\geq2$, let
$t\geq\lceil\log_2A\rceil$ be an integer, and put $B=2t+3$.
Give the symbols of $\Gamma$ distinct
labels $(a_1,\ldots,a_t)\in\{0,1\}^t$. Let $c(a)$ be the prefix
$110$ followed, in order, by the two-bit words $a_j0$ for
$j=1,\ldots,t$.
Let $\mathcal B$ extend $c$ by concatenation. For every integer
$n\geq0$ and every $x,y\in\Gamma^n$,
\begin{equation}\label{eq:delimiter}
 \frac12r_n(x,y)
 \leq\ED(\mathcal B(x),\mathcal B(y))
 \leq B\ED(x,y).
\end{equation}
\end{lemma}
\begin{proof}
The pattern $11$ occurs in a concatenation only at block starts.
The final zero of each block prevents an additional occurrence across
a boundary. Fix a longest
common subsequence matching of the encoded words, and let $l$ be its
unmatched count on each side.

At most $l$ source blocks contain an unmatched source bit. Among
the fully matched source blocks, consider one whose image is not a
consecutive target interval. Two consecutive bits of that source
block then have matched target positions with a nonempty gap. No target
position in that gap is matched: there is no source position between
the two consecutive source bits. The gaps chosen for different source
blocks are disjoint by monotonicity. Thus at most $l$ additional
source blocks have a nonconsecutive target image.

Every remaining source block is matched to a consecutive target
interval of length $B$. Its initial $11$ starts at a target block
boundary, so the entire block equals that target block. These target
blocks are distinct and occur in increasing order. Their original
letters form a common subsequence of $x,y$ with at least $n-2l$
letters. Hence $r_n(x,y)\leq2l$. By \eqref{eq:deficit}, encoded edit
distance is at least its one-sided deficit $l$, proving the lower
bound. Simulating each original symbol edit by at most $B$ bit edits
gives the upper bound, with no restriction on intermediate lengths.
When $n=0$, both sides are zero.
\end{proof}

By \eqref{eq:deficit}, the lower bound in \eqref{eq:delimiter} is at
least $\ED(x,y)/4$. Consequently, if a finite set of equal-length
words has $\ellone$ distortion at least $D_0$, its encoded image has
distortion at least $D_0/(4B)$: composing an embedding of the encoded
image with $\mathcal B$ incurs at most the factor $4B$.
